\documentclass[10pt]{article}

\usepackage[utf8]{inputenc}
\usepackage[T1]{fontenc}
\usepackage{lmodern}
\usepackage{amsmath,amssymb}
\usepackage{booktabs}
\usepackage{graphicx}
\usepackage{hyperref}
\usepackage{enumitem}
\usepackage[margin=1in]{geometry}
\usepackage{natbib}
\bibliographystyle{unsrtnat}

\title{\textbf{Supporting Information: Antipodal Complementarity in a 3$\times$3 Topological Partition of China's Climate}}

\author{Jinguo Ling}

\date{}

\begin{document}
\maketitle

This Supporting Information (SI) contains supplementary text, tables, and figures that support the main text.

\section*{SI Text}

% ══════════════════════════════════════════
\subsection*{1. Palace-Level Data}

\textbf{Table S1.} Palace-level seasonal amplitude, latitude, and antipodal pair means (UDel v501, 1961--1990 reference period, full amplitude).

\vspace{6pt}

\begin{tabular}{cccccc}
\toprule
Palace & Direction & Latitude ($^\circ$N) & Amplitude ($^\circ$C) & Antipodal pair & Pair mean ($^\circ$C) \\
\midrule
1 & N   & 48.28 & 40.32 & 1$\leftrightarrow$9 & 28.6 \\
2 & SW  & 21.93 & 10.94 & 2$\leftrightarrow$8 & 25.9 \\
3 & E   & 34.61 & 26.00 & 3$\leftrightarrow$7 & 25.0 \\
4 & SE  & 23.01 & 21.10 & 4$\leftrightarrow$6 & 28.7 \\
5 & C   & 34.60 & 26.80 & --- & --- \\
6 & NW  & 47.71 & 36.39 & & \\
7 & W   & 34.63 & 23.97 & & \\
8 & NE  & 46.59 & 40.77 & & \\
9 & S   & 21.39 & 16.85 & & \\
\bottomrule
\end{tabular}

\vspace{6pt}
\footnotesize{*Pair means are rounded to one decimal place; the reported paired CV $= 0.061$ is computed from unrounded pair means (28.585, 25.855, 24.985, 28.745).} \normalsize

\vspace{6pt}
\noindent Paired CV $= 0.061$; latitude-predicted CV $= 0.013$; $R^2$ (amp $\sim$ lat) $= 0.917$.

% ══════════════════════════════════════════
\subsection*{2. Gnomon Constraint Derivation [Historical Note]}

The Luoshu central palace number 5 encodes a gnomon constraint: for an 8-\textit{chi} gnomon, $\tan(h_\odot) = 5$ corresponds to summer solstice noon shadow of $8/5 = 1.6$ \textit{chi}. The 8-\textit{chi} gnomon is attested in the \textit{Zhoubi Suanjing}~(9) and the \textit{Huainanzi} (both Western Han), establishing it as the standard observational instrument of the period.

The latitude satisfying $\tan(90^\circ - \varphi + \varepsilon) = 5$ is $\varphi = 90^\circ + \varepsilon - \arctan(5)$. The epoch $\varepsilon \approx 23.52^\circ$ (Western Han, $\sim$100 BC) is selected because this is the period of the \textit{Zhoubi Suanjing}, the earliest text recording the ``shadow 1 \textit{chi} 6 \textit{cun}'' measurement; this choice is textually motivated rather than post hoc. This yields $\varphi \approx 34.83^\circ$N, within 0.21$^\circ$ of Luoyang.

The result is robust to epoch variation: modern $\varepsilon \approx 23.44^\circ$ yields $\varphi \approx 34.75^\circ$N (offset 0.13$^\circ$), early Shang $\varepsilon \approx 23.76^\circ$ yields $\varphi \approx 35.07^\circ$N (offset 0.45$^\circ$); Western Han and modern epochs produce latitudes within ancient observational precision ($\pm 0.5^\circ$) of Luoyang, while the early Shang value lies at the boundary (0.45$^\circ$).

While the numeral 5 admits multiple interpretations in Luoshu scholarship (Five Phases, Five Directions, Five Tones), the gnomon interpretation is uniquely testable against astronomical geography: it generates a specific, falsifiable latitude prediction that matches the historical center within measurement precision. This is independently corroborated by the \textit{Zhoubi Suanjing} recording ``summer solstice shadow 1 \textit{chi} 6 \textit{cun}'' at Yangcheng (34.4$^\circ$N), a site adjacent to Luoyang. The convergence of two independent traditions---the Luoshu numeral 5 and the \textit{Zhoubi} shadow measurement---on the same astronomical constraint establishes Luoyang as astronomically consistent rather than arbitrary.

% ══════════════════════════════════════════
\subsection*{3. Nine Provinces Mapping}

The canonical Nine Provinces mapping, documented in the \textit{Wuxing Dayi} (Sui dynasty) and traceable to the \textit{Chunqiuwei$\cdot$Wenyaogou} (Han dynasty), assigns each palace to a province: 1 (N) $\to$ Ji, 2 (SW) $\to$ Jing, 3 (E) $\to$ Qing, 4 (SE) $\to$ Xu, 5 (center) $\to$ Yu, 6 (NW) $\to$ Yong, 7 (W) $\to$ Liang, 8 (NE) $\to$ Yan, 9 (S) $\to$ Yang. The \textit{Shiji} records that the Duke of Zhou established Luoyang as ``the center of the world, where tribute routes from the four directions are equal in distance.'' The He Zun inscription (early Western Zhou, $\sim$1040 BC) records ``residing in this central state'' (\u5b85\u5179\u4e2d\u570b), the earliest known occurrence of \textit{Zhongguo}, referring to the Luoyang region.

Only 4/9 provinces achieve full directional correspondence---defined as the provincial capital falling within the corresponding azimuthal sector. The four matching positions---center (5, Yu Province = Henan/Luoyang), north (1, Ji Province = Hebei), east (3, Qing Province = Shandong), northwest (6, Yong Province = Shaanxi)---are structural key points anchoring the discrete framework. The remaining five positions are boundary positions: Jing (palace 2, SW) extends into the south-central region due to the Tibetan Plateau displacing the SW sector; Yang (palace 9, S) occupies the southeast coast rather than pure south. This structure-versus-boundary distinction supports the interpretation that Luoshu is a top-down decomposition (imposing discrete antipodal structure on geography) rather than a bottom-up classification: the structural positions are determined by the astronomical coordinate system, while the boundary positions require geographic adaptation.

% ══════════════════════════════════════════
\subsection*{4. Central Palace Anchoring}

The central palace monthly mean is 12.8$^\circ$C versus 9.9$^\circ$C for the eight outer palaces' mean. Linear regression yields $r = 0.999$ and RMSE $= 2.93^\circ$C (23\% of the central palace mean, 30\% of the outer palaces' mean), indicating effective spatial anchoring despite the systematic offset from the central palace's mid-latitude position. The offset in mean temperature reflects the central palace's mid-latitude position, which is warmer than the eight-outer-palace average that includes high-latitude (cold) and low-latitude (warm) contributions.

The central palace comprises only 63 grid points (compared to 832--2,449 for the outer palaces), limiting the precision of anchoring statistics. However, the 2.0$^\circ$ radius captures a meaningful climate unit centered on Luoyang.

% ══════════════════════════════════════════
\subsection*{5. Temporal Robustness: Full 18-Window Data}

\textbf{Table S2.} Paired CV, rank, and 3$\leftrightarrow$7 difference across 18 sliding 30-year windows.

\vspace{6pt}

\begin{tabular}{cccc}
\toprule
Window & Paired CV & Rank (of 105) & 3$\leftrightarrow$7 diff ($^\circ$C) \\
\midrule
1901--1930 & 0.0626 & 1 & 1.75 \\
1906--1935 & 0.0633 & 1 & 1.97 \\
1911--1940 & 0.0628 & 1 & 2.12 \\
1916--1945 & 0.0622 & 1 & 2.12 \\
1921--1950 & 0.0616 & 1 & 2.09 \\
1926--1955 & 0.0619 & 1 & 1.88 \\
1931--1960 & 0.0623 & 1 & 1.87 \\
1936--1965 & 0.0606 & 1 & 1.99 \\
1941--1970 & 0.0606 & 1 & 2.06 \\
1946--1975 & 0.0617 & 1 & 1.88 \\
1951--1980 & 0.0627 & 1 & 1.69 \\
1956--1985 & 0.0619 & 1 & 1.97 \\
1961--1990 & 0.0611 & 1 & 2.03 \\
1966--1995 & 0.0622 & 1 & 1.77 \\
1971--2000 & 0.0628 & 1 & 1.58 \\
1976--2005 & 0.0644 & 1 & 1.66 \\
1981--2010 & 0.0672 & 2 & 1.81 \\
1986--2015 & 0.0678 & 2 & 1.47 \\
\bottomrule
\end{tabular}

\vspace{6pt}
\noindent Key summary statistics: Luoshu ranks 1st in 16/18 windows and 2nd in 2/18 windows. Meta-CV $= 0.030$. Mean 3$\leftrightarrow$7 difference $= 1.87^\circ$C (positive in all windows). The latitude degradation direction (observed CV $>$ latitude-predicted CV) is consistent across all windows.

Weak trends are detected: CV upward trend (slope $= +0.000037$/yr, $p = 0.031$) and 3$\leftrightarrow$7 declining trend (slope $= -0.0045^\circ$C/yr, $p = 0.006$). Both are consistent with global warming signals: Arctic amplification reduces the north--south amplitude gradient (increasing CV), while faster warming in the elevated northwest reduces the east--west amplitude difference (decreasing 3$\leftrightarrow$7). Neither trend alters qualitative conclusions.

% ══════════════════════════════════════════
\subsection*{6. EOF2 Sign-Reversal Test: Methodological Caveat}

A natural diagnostic for antipodal coupling is checking whether EOF patterns exhibit sign reversal between antipodal sectors. This test is insufficient as a diagnostic for antipodal coupling in eight azimuthal sectors under continuous spatial gradients, for two independent reasons:

\textbf{General principle: sign is mathematically arbitrary.} Eigenvector signs are indeterminate---multiplication by $-1$ preserves eigenvalue and eigenvector status~(11). North \textit{et al.}~(11) further showed that closely spaced eigenvalues can cause eigenvector mixing---a related but distinct issue concerning modal separation rather than sign direction.

\textbf{Specific condition: geometric inevitability.} When eight equal-width azimuthal sectors span a continuous meridional gradient, antipodal pairs inevitably lie at opposite gradient ends, producing sign reversal as geometric necessity. A random-rotation null model yields P(4/4 sign reversal) $= 0.412$ (non-significant).

The two layers reinforce each other: even setting aside the general arbitrariness of EOF signs, the specific geometry of azimuthal sectors guarantees sign reversal under any monotonic gradient. Sign reversal therefore cannot differentiate Luoshu-specific antipodal coupling from geometric triviality. By contrast, amplitude complementarity (CV) depends on the specific pairing rule and can differentiate structures---making it a valid diagnostic. This distinction---valid (CV complementarity) versus insufficient (sign reversal)---is a methodological contribution applicable to any spatial framework that partitions a continuous field into discrete sectors.

% ══════════════════════════════════════════
\subsection*{7. Cross-Dataset Validation}

The main analysis uses UDel v501 (0.5$^\circ$, 1900--2017). To assess dataset dependence, we replicate the palace-level amplitude analysis using the CUG-CMA CHDT dataset~(20) (2.5$^\circ$ resolution, 1880--2020; reference period 1961--1990). Despite the different spatial resolution and temporal coverage, the CUG-CMA results are consistent with UDel: the Luoshu pairing achieves rank 1/105 (CV $= 0.066$ vs.\ UDel 0.061), and the inter-annual amplitude series are strongly correlated (Spearman $\rho = 0.712$, $p < 0.0001$). Both datasets independently place the Luoshu pairing at the top rank, confirming that the result is not an artifact of a specific gridded dataset.

% ══════════════════════════════════════════
\subsection*{8. Robustness Checks}

\textbf{Sensitivity to central point.} Five candidate cities tested: Luoyang, Xi'an, Kaifeng, Jingyang, Beijing. Luoyang yields the lowest paired CV; all centers achieve $p < 0.05$ in Tier A and Tier B.

\textbf{Sensitivity to central palace radius.} Radius varied from 1$^\circ$ to 5$^\circ$; CV and correlation remain stable across this range.

\textbf{Equal-weighted vs.\ area-weighted amplitudes.} The nine palaces vary substantially in area: the largest outer palace (7, W) has 2,449 grid points versus the smallest (1, N) with 832 points (2.9$\times$ ratio). Area-weighting increases rather than decreases the 3$\leftrightarrow$7 contrast (paired CV: 0.0611 vs.\ 0.0628, 2.7\% difference; 3$\leftrightarrow$7: 2.03$^\circ$C vs.\ 2.49$^\circ$C area-weighted), falsifying the hypothesis that area inequality artificially inflates the east--west signal.

\textbf{Linear vs.\ quadratic regression.} With only eight data points, a quadratic model ($\text{amp} \sim \text{lat} + \text{lat}^2$) yields $R^2 = 0.924$ ($\Delta R^2 = 0.006$) at the cost of an additional parameter; the adjusted $R^2$ decreases from 0.904 (linear) to 0.893 (quadratic), confirming overfitting~(13). The quadratic model also introduces a physically implausible curvature in the latitude--amplitude relationship over the observed range.

% ══════════════════════════════════════════
\subsection*{8. Comparison with Data-Driven Climate Regionalization}

A natural question is how the Luoshu framework compares with data-driven regionalization methods such as $k$-means clustering, hierarchical clustering, or rotated EOF analysis~(6, 12, 13). The comparison must account for parameter efficiency: data-driven methods with many free parameters will trivially achieve better fit, just as a polynomial of sufficient degree can interpolate any set of points. Luoshu specifies 9 regions with effectively 2 degrees of freedom (the center point and central palace radius; sector angles and pairing structure are fixed by the Luoshu rule), whereas $k$-means with $k = 9$ has 18 free parameters (2D centroids).

The fair comparison is not fit (where data-driven methods win by construction) but generalizability: does the prespecified structure capture detectable organization in data it was not designed to fit? The exhaustive enumeration (Tier C) provides this test: Luoshu achieves rank 1/105 among all possible pairings, demonstrating that its structure is not merely compatible with the data but optimal---despite having no free parameters in the pairing rule.

Buskop \textit{et al.}~(14) employed sub-catchment partitioning in compound-event analysis; we extend this discrete-partitioning principle from physically defined regions (catchments) to culturally encoded regions (Luoshu palaces), demonstrating that the latter carries detectable statistical structure.

% ══════════════════════════════════════════
\subsection*{9. Monthly Complementarity}

\textbf{Table S3.} Luoshu pairing rank among 105 schemes for each calendar month (NCEP/NCAR reanalysis surface air temperature, 1961--1990 reference period).

\vspace{6pt}

\begin{tabular}{ccc}
\toprule
Month & Rank (of 105) & Note \\
\midrule
Jan & 5 & Strong N--S gradient \\
Feb & 2 & \\
Mar & 1 & \\
Apr & 1 & \\
May & 1 & \\
Jun & 1 & Weakest gradient \\
Jul & 2 & \\
Aug & 2 & \\
Sep & 4 & \\
Oct & 4 & \\
Nov & 6 & \\
Dec & 6 & Strong N--S gradient \\
\bottomrule
\end{tabular}

\vspace{6pt}
\noindent All 12 months achieve rank $\leq$6/105 (mean $= 2.9$). We evaluate monthly generality by rank rather than by the absolute CV value, for the following reason. The monthly analysis uses palace-mean absolute temperature (not seasonal amplitude); when the mean of the four pair means approaches 0$^\circ$C (as occurs in transitional and winter months), CV $= \sigma/\mu$ becomes unstable. However, each of the eight outer palaces appears exactly once in every pairing scheme, so the sum $S = \sum_i T_i$ is identical across all 105 schemes, making $\mu = S/8$ of the four pair means a constant independent of the pairing rule. The CV ranking is therefore equivalent to the $\sigma$ ranking and remains valid. Summer months (May--July) rank 1--2, reflecting a weaker meridional temperature gradient that makes alternative pairings less competitive. Winter months (November--January) rank 5--6, reflecting the stronger gradient. The seasonal modulation of rank is consistent with Contribution~1 (latitude gradient): when the latitude gradient is strong, more pairing schemes can partially exploit it, increasing competition for the top rank; when the gradient is weak, Luoshu's antipodal structure becomes more uniquely advantageous.

% ══════════════════════════════════════════
\subsection*{10. Multi-Variable Complementarity}

\textbf{Table S4.} Luoshu pairing rank among 105 schemes for five NCEP/NCAR reanalysis variables (1961--1990 reference period, seasonal amplitude metric).

\vspace{6pt}

\begin{tabular}{lcccc}
\toprule
Variable & Luoshu CV & Rank (of 105) & $r$ (lat--amp) & $p$ (lat--amp) \\
\midrule
TMAX & 0.081 & 1 & 0.989 & $<0.001$ \\
TMEAN & 0.090 & 1 & 0.989 & $<0.001$ \\
TMIN & 0.108 & 8 & 0.971 & $<0.001$ \\
Wind speed & 0.270 & 44 & $-0.449$ & 0.225 \\
Total cloud cover & 0.313 & 74 & $-0.643$ & 0.085 \\
\bottomrule
\end{tabular}

\vspace{6pt}
\noindent The gradient from rank 1 (TMAX, TMEAN) through rank 8 (TMIN) to rank 74 (cloud cover) tracks the latitude--amplitude correlation: temperature variables have strongly positive $r$ (amplitude increases with latitude, driven by solar radiation), wind speed has weakly negative $r$ (no systematic latitude dependence), and cloud cover has moderately negative $r$ (cloud fraction decreases with latitude over China). This gradient supports Contribution~1: the latitude gradient is the primary physical driver of antipodal complementarity, and variables whose seasonal amplitude is less latitude-dependent exhibit correspondingly weaker specificity.

The TMIN rank (8) being higher than TMAX (1) reflects the greater complexity of minimum temperature, which is more strongly influenced by local factors (cloud cover, humidity, urban heat islands) than maximum temperature, which is primarily radiation-driven. The wind speed intermediate rank (44) reflects the mixed signal: while wind speed amplitude has some latitude structure, it is dominated by the monsoon--topography interaction that varies across the nine palaces in ways not aligned with the antipodal structure.

\section*{SI References}

\begin{thebibliography}{19}
\bibitem{6} H.~E.~Beck \textit{et al.}, \textit{Sci.~Data} \textbf{10}, 724 (2023).
\bibitem{9} C.~Cullen, \textit{Astronomy and Mathematics in Ancient China: The Zhou Bi Suan Jing} (Cambridge University Press, 1996).
\bibitem{11} G.~R.~North \textit{et al.}, \textit{Mon.~Weather Rev.} \textbf{110}, 699--706 (1982).
\bibitem{12} H.~Duan \textit{et al.}, \textit{Atmosphere} \textbf{15}, 741 (2024).
\bibitem{13} H.~Akaike, \textit{IEEE Trans.~Autom.~Control} \textbf{19}, 716--723 (1974).
\bibitem{14} F.~E.~Buskop \textit{et al.}, \textit{Earth's Future} \textbf{13}, e2025EF006499 (2025).
\bibitem{19} E.~Kalnay \textit{et al.}, \textit{Bull.~Am.~Meteorol.~Soc.} \textbf{77}, 437--471 (1996).
\bibitem{20} X.~Zheng, Y.~Ren, G.~Ren, P.~Zhang \textit{et al.}, A century-long China homogenized daily surface air temperature dataset (CUG-CMA CHDT). \textit{Sci.~Data} \textbf{11}, 1045 (2024). DOI: 10.1038/s41597-024-03758-3.
\end{thebibliography}

\end{document}